\clearpage
\VerdierPhysicalPageAnchor{0058}%
\VerdierSemanticAnchor{ch1-1-5-10-converse}
\noindent\textbf{反过来：}\par

\VerdierSemanticAnchor{ch1-1-5-10}
\VerdierSemanticAnchor{ch1-1-5-10-proposition}
\noindent\textbf{命题 1.5.10.}\quad
\emph{设 $u\colon F_1\to F_2$ 为底层函子之间的一个态射；若对每个
$i\in[1,n]$，图 (1.5.9.2) 都为 $\beta(e_i)$-交换的，
则 $u$ 是函子的一个 $\beta$-态射。}

\VerdierSemanticAnchor{ch1-1-5-10-proof}
事实上，对每个 $\sigma\in\mathbb Z^n$，以 $R(\sigma)$ 表示性质：

\VerdierSemanticAnchor{ch1-1-5-10-property}
\noindent\emph{$(R(\sigma))$\qquad 下图为 $\beta(\sigma)$-交换的：}

\VerdierSemanticAnchor{ch1-1-5-10-eq1}
\begin{equation}
  \resizebox{0.78\textwidth}{!}{%
    $\begin{tikzcd}[ampersand replacement=\&,column sep=6em,row sep=4em]
      F_1T(\sigma)
        \arrow[r,"{p_1(\sigma)}"]
        \arrow[d,"{u*T(\sigma)}"']
        \& T'(\sigma)F_1
          \arrow[d,"{T'(\sigma)*u}"] \\
      F_2T(\sigma)
        \arrow[r,"{p_2(\sigma)}"]
        \& T'(\sigma)F_2
    \end{tikzcd}$%
  }\;.
  \tag*{\ensuremath{(1.5.10.1)_{\sigma}}}
\end{equation}

考察各图 $(1.5.10.1)_{\sigma_1}$、$(1.5.10.1)_{\sigma_2}$ 与
$(1.5.10.1)_{\sigma_1+\sigma_2}$，立即得到：

\VerdierSemanticAnchor{ch1-1-5-10-closure}
\noindent\emph{若性质 $R(\sigma_1)$、$R(\sigma_2)$、
$R(\sigma_1+\sigma_2)$ 中有两个成立，则第三个也成立。}

\VerdierSemanticAnchor{ch1-1-5-10-conclusion}
此外，由假设，性质 $R(e_i)$（$i\in[1,n]$）与 $R(0)$ 均成立。
由此立即可知，$R(\sigma)$ 对每个 $\sigma\in\mathbb Z^n$ 都成立。

\VerdierSemanticAnchor{ch1-1-5-11}
\noindent\textbf{1.5.11.}\quad
一个严格 $\mathbb Z^n$-范畴，是一个装备了 $\mathbb Z^n$ 到其自同构群
的表示的范畴。要定义这样的表示，当然只需给出 $n$ 个两两交换的自同构
$T(e_1),\ldots,T(e_n)$。

\VerdierSemanticAnchor{ch1-1-5-12}
\noindent\textbf{1.5.12.}\quad
在一个给定范畴上确定（不必严格的）$\mathbb Z^n$-范畴结构，问题要复杂得多。
例如，只考虑范畴 $\mathcal C$ 上满足每个 $\sigma\in\mathbb Z^n$ 都有
$T(\sigma)=\text{$\mathcal C$ 上的恒等函子}$ 的 $\mathbb Z^n$-范畴结构。
这时必须对每一对 $\sigma_1,\sigma_2\in\mathbb Z^n$，给出
$\mathcal C$ 的恒等函子的自同构群中的一个元素
$c(\sigma_1,\sigma_2)$（该群为交换群），并要求这些数据满足恒等式
(1.2.1.3)。因此，要给出这样的结构，就必须给出一个取值于
$\mathcal C$ 的恒等函子自同构群的、$\mathbb Z^n$ 上的正规化非齐次
2-余循环。还可立即验证：两个这样的结构可由一个满足
$F=\text{$\mathcal C$ 上的恒等函子}$ 的 $\mathbb Z^n$-函子 $(F,m)$
同构，当且仅当相应余循环属于同一上同调类。
